% Part: incompleteness
% Chapter: incompleteness-provability
% Section: lob-thm

\documentclass[../../../include/open-logic-section]{subfiles}

\begin{document}

\olfileid{inc}{inp}{lob}

\olsection{قضیهٔ لوب}

جملهٔ گودلیِ نظریهٔ~$\Th{T}$ نقطهٔ ثابتی برای~$\lnot
\OProv[\Th{T}](y)$ است، یعنی !!a{sentence} چون~$!G$ که
\[
\Th{T} \Proves \lnot \OProv[\Th{T}](\gn{!G}) \liff !G.
\]
اگر~$\Th{T}$ سازگار باشد، این جمله !!{derivable} نیست؛ زیرا اگر
$\Th{T} \Proves !G$، آنگاه (الف) بنا بر شرط~(۱) !!{derivability}،
$\Th{T} \Proves \OProv[\Th{T}](\gn{!G})$، و (ب) از
$\Th{T} \Proves !G$ همراه با
$\Th{T} \Proves \lnot \OProv[\Th{T}](\gn{!G}) \liff !G$ به دست
می‌آید~$\Th{T} \Proves \lnot \OProv[\Th{T}](\gn{!G})$؛ بنابراین
$\Th{T}$ ناسازگار می‌شد. اکنون طبیعی است دربارهٔ وضعیت نقطهٔ ثابتی
برای~$\OProv[\Th{T}](y)$ بپرسیم، یعنی !!a{sentence} چون~$!H$ که
\[
\Th{T} \Proves \OProv[\Th{T}](\gn{!H}) \liff !H.
\]
اگر این جمله !!{derivable} بود، بنا بر شرط~(۱) داشتیم
$\Th{T} \Proves \OProv[\Th{T}](\gn{!H})$؛ اما همین نتیجه با اعمال
قیاس استثنایی بر هم‌ارزی بالا نیز به دست می‌آید. پس، دست‌کم نه با همان
استدلالی که در مورد جملهٔ گودلی به کار رفت، به ناسازگاری~$\Th{T}$
نمی‌رسیم. البته این نشان نمی‌دهد که~$\Th{T}$ \emph{واقعاً} $!H$ را
!!{derive} می‌کند.

اگر پرسش را اندکی تعمیم دهیم می‌توانیم در پاسخ به آن پیشرفت کنیم. جهت
چپ‌به‌راستِ هم‌ارزی نقطهٔ ثابت،
$\OProv[\Th{T}](\gn{!H}) \lif !H$، نمونه‌ای از شمایی کلی است که
\emph{اصل بازتاب} نامیده می‌شود:
$\OProv[\Th{T}](\gn{!A}) \lif !A$. این نام از آن روست که شِما، به یک
معنا، بیان می‌کند~$\Th{T}$ می‌تواند دربارهٔ آنچه !!{derive} می‌کند
«بازتاب» داشته باشد؛ اساساً می‌گوید: «اگر~$\Th{T}$ بتواند~$!A$ را
!!{derive} کند، آنگاه~$!A$ صادق است»، برای هر~$!A$. البته این فقط
دربارهٔ نظریه‌های درست صادق است و نشان می‌دهد نظریه‌ها عموماً هر نمونهٔ
آن را !!{derive} نمی‌کنند. پس یک نظریه (که به‌اندازهٔ کافی قوی باشد و
شرایط !!{derivability} را برآورده کند) کدام نمونه‌ها را می‌تواند
!!{derive} کند؟ بی‌گمان همهٔ نمونه‌هایی که خود~$!A$ در آن‌ها
!!{derivable} است. همین و بس، چنان‌که نتیجهٔ بعد نشان می‌دهد.

\begin{thm}\ollabel{thm:lob}
فرض کنید~$\Th{T}$ نظریه‌ای !!a{axiomatizable} و گسترشی از~$\Th{Q}$ باشد
و~$\OProv[\Th{T}](y)$ فرمولی باشد که شرایط P1--P3 از~\olref[2in]{sec}
را برآورده می‌کند. اگر~$!A$ یک جمله باشد و~$\Th{T}$،
$\OProv[\Th{T}](\gn{!A}) \lif !A$ را !!{derive}s کند، آنگاه در واقع
$\Th{T}$، $!A$ را !!{derive}s می‌کند.
\end{thm}

به بیان دیگر، اگر~$\Th{T} \Proves/ !A$، آنگاه
$\Th{T} \Proves/ \OProv[\Th{T}](\gn{!A}) \lif !A$. این نتیجه به
قضیهٔ لوب معروف است.

\begin{explain}
شهود راهنمای اثبات قضیهٔ لوب، اثباتی هوشمندانه برای وجود بابانوئل است.
(اگر این نتیجه را نمی‌پسندید، می‌توانید هر نتیجهٔ دیگری را که می‌خواهید
جایگزین کنید.) استدلال چنین است:
\begin{enumerate}
\item بگذارید~$X$ جملهٔ «اگر~$X$ صادق باشد، آنگاه بابانوئل وجود دارد»
  باشد.
\item فرض کنید~$X$ صادق باشد.
\item آنگاه آنچه می‌گوید برقرار است؛ یعنی داریم: اگر~$X$ صادق باشد،
  آنگاه بابانوئل وجود دارد.
\item چون فرض می‌کنیم~$X$ صادق است، با قیاس استثنایی از (۲) و~(۳)
  می‌توانیم نتیجه بگیریم بابانوئل وجود دارد.
\item توانستیم (۴)، یعنی «بابانوئل وجود دارد»، را از فرض~(۲)، یعنی
  «$X$ صادق است»، به دست آوریم. با اثبات شرطی نشان داده‌ایم: «اگر~$X$
  صادق باشد، آنگاه بابانوئل وجود دارد.»
\item اما این دقیقاً همان جملهٔ~$X$ است. پس نشان داده‌ایم~$X$ صادق است.
\item اما آنگاه بنا بر استدلال (۲)--(۴) بالا، بابانوئل وجود دارد.
\end{enumerate}
صوری‌سازی این ایده، با جایگزینی «صادق است» با «!!{derivable} است» و
«بابانوئل وجود دارد» با~$!A$، اثبات قضیهٔ لوب را به دست می‌دهد. ترفند
آن است که لم نقطهٔ ثابت را بر !!{formula}~$\OProv[\Th{T}](y) \lif !A$
اعمال کنیم. نقطهٔ ثابت آن متناظر با جملهٔ~$X$ در طرح پیشین است.
\end{explain}

\begin{proof}[اثبات \olref{thm:lob}]
فرض کنید~$!A$ !!a{sentence}ی باشد چنان‌که~$\Th{T}$،
$\OProv[\Th{T}](\gn{!A}) \lif !A$ را !!{derive}s می‌کند. بگذارید
$!B(y)$ همان !!{formula}~$\OProv[\Th{T}](y) \lif !A$ باشد و با لم
نقطهٔ ثابت !!a{sentence} چون~$!D$ بیابید به‌گونه‌ای که~$\Th{T}$،
$!D \liff !B(\gn{!D})$ را !!{derive}s می‌کند. آنگاه هر یک از عبارت‌های
زیر در~$\Th{T}$ !!{derivable} است:
\begin{align}
  & !D \liff (\OProv[\Th{T}](\gn{!D}) \lif !A) \ollabel{L-1}\\
  & \qquad \text{$!D$ نقطهٔ ثابتی برای~$!B(y)$ است}\notag \\
  & !D \lif (\OProv[\Th{T}](\gn{!D}) \lif !A) \ollabel{L-2}\\
  & \qquad\text{از \olref{L-1}}\notag\\
  & \OProv[\Th{T}](\gn{!D \lif (\OProv[\Th{T}](\gn{!D}) \lif !A)}) \ollabel{L-3}\\
  & \qquad \text{از \olref{L-2} بنا بر شرط P1}\notag \\
  & \OProv[\Th{T}](\gn{!D}) \lif \OProv[\Th{T}](\gn{\OProv[\Th{T}](\gn{!D}) \lif !A})
  \ollabel{L-4}\\
  &\qquad \text{از \olref{L-3} با استفاده از شرط P2}\notag \\
  & \OProv[\Th{T}](\gn{!D}) \lif (\OProv[\Th{T}](\gn{\OProv[\Th{T}](\gn{!D})}) \lif \OProv[\Th{T}](\gn{!A})) \ollabel{L-5}\\
  &\qquad \text{از \olref{L-4} با استفادهٔ دوباره از P2} \notag\\
& \OProv[\Th{T}](\gn{!D}) \lif \OProv[\Th{T}](\gn{\OProv[\Th{T}](\gn{!D})}) \ollabel{L-6}\\
  & \qquad\text{بنا بر شرط P3 از !!{derivability}} \notag\\
  & \OProv[\Th{T}](\gn{!D}) \lif \OProv[\Th{T}](\gn{!A}) \ollabel{L-7} \\
  &\qquad\text{از \olref{L-5} و \olref{L-6}}\notag\\
  & \OProv[\Th{T}](\gn{!A}) \lif !A \ollabel{L-8}\\
  &\qquad\text{بنا بر فرض قضیه} \notag\\
  & \OProv[\Th{T}](\gn{!D}) \lif !A \ollabel{L-9}\\
  &\qquad\text{از \olref{L-7} و \olref{L-8}}\notag\\
  & (\OProv[\Th{T}](\gn{!D}) \lif !A) \lif !D \ollabel{L-10}\\
  & \qquad \text{از \olref{L-1}}\notag \\
  & !D \ollabel{L-11}\\
  & \qquad\text{از \olref{L-9} و \olref{L-10}}\notag \\
  & \OProv[\Th{T}](\gn{!D}) \ollabel{L-12}\\
  & \qquad\text{از \olref{L-11} بنا بر شرط~P1}\notag \\
  & !A \qquad\qquad\text{از \olref{L-8} و \olref{L-12}}\notag
\end{align}
\end{proof}

با داشتن قضیهٔ لوب، اثبات کوتاهی برای قضیهٔ دوم ناتمامیت (برای
نظریه‌هایی با !!a{derivability} محمولی که شرایط P1--P3 را برآورده
می‌کند) داریم: اگر~$\Th{T} \Proves
\OProv[\Th{T}](\gn{\lfalse}) \lif \lfalse$، آنگاه
$\Th{T} \Proves \lfalse$. اگر~$\Th{T}$ سازگار باشد،
$\Th{T} \Proves/ \lfalse$. پس
$\Th{T} \Proves/ \OProv[\Th{T}](\gn{\lfalse}) \lif \lfalse$، یعنی
$\Th{T} \Proves/ \OCon[\Th{T}]$. همچنین می‌توانیم آن را به کار بریم
تا نشان دهیم~$!H$، یعنی نقطهٔ ثابت~$\OProv[\Th{T}](x)$، !!{derivable}
است. زیرا
\begin{align*}
  \Th{T} & \Proves \OProv[\Th{T}](\gn{!H}) \liff !H\\
  \intertext{به‌ویژه}
    \Th{T} & \Proves \OProv[\Th{T}](\gn{!H}) \lif !H
\end{align*}
و بنابراین بنا بر قضیهٔ لوب، $\Th{T} \Proves !H$.

% Going in the other direction, for homework I
% may ask you to work through a short proof of L"ob's theorem, using
% the second incompleteness theorem instead of the fixed-point lemma.

\begin{prob}
فرض کنید~$\Th{T}$ نظریه‌ای به‌طور محاسبه‌پذیر اصل‌بندی‌شده و
$\OProv[\Th{T}]$ !!a{derivability} محمولی برای~$\Th{T}$ باشد. چهار
گزارهٔ زیر را در نظر بگیرید:
\begin{enumerate}
\item اگر~$T \Proves !A$، آنگاه~$T \Proves \OProv[\Th{T}](\gn{!A})$.
\item $T \Proves !A \lif \OProv[\Th{T}](\gn{!A})$.
\item اگر~$T \Proves \OProv[\Th{T}](\gn{!A})$، آنگاه~$T \Proves !A$.
\item $T \Proves \OProv[\Th{T}](\gn{!A}) \lif !A$
\end{enumerate}
هر یک از این گزاره‌ها تحت چه شرایطی صادق است؟
\end{prob}

\end{document}
